\documentclass[12pt,a4paper]{ctexart}
\usepackage{amsmath,amssymb}
\usepackage{geometry}
\usepackage{listings}
\usepackage{xcolor}
\usepackage{hyperref}
\hypersetup{
	colorlinks=true,
	linkcolor={blue!60!black},
	urlcolor={blue!60!black},
	citecolor={blue!60!black},
	pdfborder={0 0 0},
}
\usepackage{tabularx}
\usepackage{fancyhdr}
\usepackage{tikz}
\usepackage{lastpage}
\usepackage{enumitem}
\usepackage{array}
\usepackage{booktabs}
\usepackage{longtable}
\geometry{left=2.5cm,right=2.5cm,top=2.5cm,bottom=2.5cm}
\definecolor{codebg}{RGB}{245,245,245}
\definecolor{codeframe}{RGB}{200,200,200}
\definecolor{commentcolor}{RGB}{0,128,0}
\definecolor{keywordcolor}{RGB}{127,0,85}
\definecolor{stringcolor}{RGB}{42,0,255}
\definecolor{accent}{HTML}{4F46E5}
\definecolor{accent2}{HTML}{7C3AED}
\definecolor{mid}{HTML}{6B7280}
\lstset{
	language=C++,
	basicstyle=\small\ttfamily,
	backgroundcolor=\color{codebg},
	frame=single,
	rulecolor=\color{codeframe},
	breaklines=true,
	showstringspaces=false,
	commentstyle=\color{commentcolor},
	keywordstyle=\color{keywordcolor}\bfseries,
	stringstyle=\color{stringcolor},
	numbers=left,
	numberstyle=\tiny\color{gray},
	tabsize=4,
	captionpos=b,
	extendedchars=true,
	columns=flexible,
}
\pagestyle{fancy}
\fancyhf{}
\fancyhead[L]{\leftmark}
\fancyhead[R]{数据结构与算法专题}
\fancyfoot[C]{\thepage\ /\ \pageref*{LastPage}}
\renewcommand{\headrulewidth}{0.4pt}
\renewcommand{\footrulewidth}{0pt}
\newcommand{\infobox}[1]{%
	\par\noindent
	\fcolorbox{accent!40}{accent!5}{%
		\begin{minipage}{\dimexpr\textwidth-2\fboxsep-2\fboxrule\relax}
			#1
		\end{minipage}%
	}\par
}
\title{数据结构与算法专题}
\author{算法学习笔记}
\date{\today}
\begin{document}
	\maketitle
	\tableofcontents
	\newpage
	\section{初级数据结构}
	\subsection{栈——括号匹配}
	\infobox{解决问题：验证括号序列是否合法，即每个左括号是否都有对应的右括号且类型匹配。适用于编译原理语法检查、XML/HTML标签匹配等。}
	\subsubsection*{题目描述}
	给定一个仅包含字符 '('、')'、'['、']'、'{'、'}' 的字符串 $s$，判断括号序列是否合法。合法条件：每个左括号必须由相同类型的右括号闭合，且括号嵌套顺序正确。
	\subsubsection*{输入示例}
	\begin{verbatim}
		{[]()}
	\end{verbatim}
	\subsubsection*{输出示例}
	YES
	\begin{lstlisting}
		#include <bits/stdc++.h>
		using namespace std;
		int main() {
			string s;               // s: 输入的括号字符串
			cin >> s;
			stack<char> st;         // st: 栈，用来存放未匹配的左括号
			bool ok = true;         // ok: 记录括号序列是否合法
			for (char c : s) {      // c: 当前字符
				if (c == '(' || c == '[' || c == '{') {
						st.push(c);     // 左括号直接压栈
					} else {
						if (st.empty()) { ok = false; break; }  // 栈空，无匹配左括号
						char top = st.top();                    // top: 栈顶左括号
						if ((c == ')' && top != '(') ||
						(c == ']' && top != '[') ||
						(c == '}' && top != '{')) {
							ok = false; break;                  // 类型不匹配，非法
						}
						st.pop();       // 匹配成功，弹出左括号
					}
				}
				if (!st.empty()) ok = false;   // 最后栈不空说明有左括号未被匹配
				cout << (ok ? "YES" : "NO") << "\n";
				return 0;
			}
		\end{lstlisting}
		\subsection{栈——单调栈求下一个更大元素}
		\infobox{解决问题：对数组中每个元素，找到其右侧第一个比它大的元素。适用于股票价格跨度、每日温度等待天数等场景。}
		\subsubsection*{题目描述}
		给定长度为 $n$ 的数组 $a$，对于每个位置 $i$，找出 $a[i]$ 右侧第一个大于 $a[i]$ 的元素；若不存在则输出 $-1$。
		\subsubsection*{输入示例}
		\begin{verbatim}
			5
			2 1 3 5 4
		\end{verbatim}
		\subsubsection*{输出示例}
		3 3 5 -1 -1
		\begin{lstlisting}
			#include <bits/stdc++.h>
			using namespace std;
			int main() {
				int n;                    // n: 数组长度
				cin >> n;
				vector<int> a(n);         // a: 输入数组
				for (int i = 0; i < n; i++) cin >> a[i];
				vector<int> ans(n, -1);   // ans[i]: a[i]右侧第一个更大的元素，初始为-1
				stack<int> st;            // st: 单调栈，存放下标，栈底到栈顶对应值递减
				for (int i = 0; i < n; i++) {   // i: 当前元素下标
					while (!st.empty() && a[st.top()] < a[i]) {
						ans[st.top()] = a[i];   // 当前元素是栈顶元素右侧第一个更大的
						st.pop();
					}
					st.push(i);
				}
				for (int i = 0; i < n; i++) cout << ans[i] << " \n"[i==n-1];
				return 0;
			}
		\end{lstlisting}
		\subsection{队列——滑动窗口最大值}
		\infobox{解决问题：在长度为 $n$ 的数组中，求所有长度为 $k$ 的连续子数组的最大值。适用于数据流分析、信号处理等滑动窗口场景。}
		\subsubsection*{题目描述}
		给定数组 $a$ 和窗口大小 $k$，输出每个窗口内的最大值。
		\subsubsection*{输入示例}
		\begin{verbatim}
			8 3
			1 3 -1 -3 5 3 6 7
		\end{verbatim}
		\subsubsection*{输出示例}
		3 3 5 5 6 7
		\begin{lstlisting}
			#include <bits/stdc++.h>
			using namespace std;
			int main() {
				int n, k;                 // n: 数组长度, k: 窗口大小
				cin >> n >> k;
				vector<int> a(n);         // a: 输入数组
				for (int i = 0; i < n; i++) cin >> a[i];
				deque<int> dq;            // dq: 双向队列，存放下标，队首为当前窗口最大值下标
				for (int i = 0; i < n; i++) {
					while (!dq.empty() && dq.front() <= i - k) dq.pop_front(); // 移除超出窗口的下标
					while (!dq.empty() && a[dq.back()] <= a[i]) dq.pop_back(); // 保持队列递减
					dq.push_back(i);
					if (i >= k - 1) cout << a[dq.front()] << " "; // 窗口形成后输出最大值
				}
				return 0;
			}
		\end{lstlisting}
		\subsection{并查集——连通分量与动态合并}
		\infobox{解决问题：动态维护若干集合的合并与查询，判断两个元素是否属于同一集合。适用于社交网络好友关系、岛屿连通性、Kruskal最小生成树等。}
		\subsubsection*{题目描述}
		给定 $n$ 个节点和 $m$ 条无向边，再给出 $q$ 个询问，每次询问两个节点是否连通。
		\subsubsection*{输入示例}
		\begin{verbatim}
			5 3
			1 2
			2 3
			4 5
			2
			1 3
			4 5
		\end{verbatim}
		\subsubsection*{输出示例}
		YES YES
		\begin{lstlisting}
			#include <bits/stdc++.h>
			using namespace std;
			vector<int> parent, rnk;    // parent[i]: i的父节点; rnk[i]: 以i为根的秩（近似树高）
			int find(int x) {           // 查找x的根，同时进行路径压缩
				return parent[x] == x ? x : parent[x] = find(parent[x]);
			}
			void unite(int a, int b) {  // 合并a和b所在的集合（按秩合并）
				int ra = find(a), rb = find(b);
				if (ra == rb) return;
				if (rnk[ra] < rnk[rb]) swap(ra, rb);
				parent[rb] = ra;
				if (rnk[ra] == rnk[rb]) rnk[ra]++;
			}
			int main() {
				int n, m;                 // n: 节点数, m: 边数
				cin >> n >> m;
				parent.resize(n+1); rnk.resize(n+1, 0);
				for (int i = 1; i <= n; i++) parent[i] = i;
				for (int i = 0; i < m; i++) {
					int u, v;             // u,v: 边的两个端点
					cin >> u >> v;
					unite(u, v);
				}
				int q;                    // q: 询问数量
				cin >> q;
				while (q--) {
					int a, b;             // a,b: 询问的两个节点
					cin >> a >> b;
					cout << (find(a) == find(b) ? "YES" : "NO") << "\n";
				}
				return 0;
			}
		\end{lstlisting}
		\subsection{哈希表——两数之和}
		\infobox{解决问题：在数组中找出两个数，使它们的和等于目标值。适用于金融对账、组合优化等查找配对场景。}
		\subsubsection*{题目描述}
		给定整数数组 $a$ 和目标值 $target$，找出数组中两个不同位置的数，使得它们的和为 $target$，输出它们的下标（从0开始）。
		\subsubsection*{输入示例}
		\begin{verbatim}
			4 9
			2 7 11 15
		\end{verbatim}
		\subsubsection*{输出示例}
		0 1
		\begin{lstlisting}
			#include <bits/stdc++.h>
			using namespace std;
			int main() {
				int n, target;            // n: 数组长度, target: 目标和
				cin >> n >> target;
				vector<int> a(n);         // a: 输入数组
				for (int i = 0; i < n; i++) cin >> a[i];
				unordered_map<int,int> mp;// mp: 值到下标的映射
				for (int i = 0; i < n; i++) {
					int need = target - a[i]; // need: 需要配对的另一个数
					if (mp.count(need)) {
						cout << mp[need] << " " << i << "\n";
						return 0;
					}
					mp[a[i]] = i;
				}
				return 0;
			}
		\end{lstlisting}
		\subsection{堆——动态中位数}
		\infobox{解决问题：在一个不断加入数字的数据流中，随时查询当前所有数字的中位数。适用于在线数据分析、实时统计等场景。}
		\subsubsection*{题目描述}
		有 $n$ 个数依次输入，每输入奇数个数时，输出当前已输入数字的中位数。
		\subsubsection*{输入示例}
		\begin{verbatim}
			7
			1 3 5 2 4 6 7
		\end{verbatim}
		\subsubsection*{输出示例}
		1 3 3
		\begin{lstlisting}
			#include <bits/stdc++.h>
			using namespace std;
			int main() {
				int n;                    // n: 数字个数
				cin >> n;
				priority_queue<int> maxHeap;                     // maxHeap: 最大堆，存较小一半
				priority_queue<int, vector<int>, greater<int>> minHeap; // minHeap: 最小堆，存较大一半
				for (int i = 1; i <= n; i++) {
					int x;                // x: 当前输入数字
					cin >> x;
					maxHeap.push(x);      // 先放入最大堆
					// 保持平衡：最大堆顶 <= 最小堆顶
					if (!minHeap.empty() && maxHeap.top() > minHeap.top()) {
						int a = maxHeap.top(); maxHeap.pop();
						int b = minHeap.top(); minHeap.pop();
						maxHeap.push(b); minHeap.push(a);
					}
					// 大小平衡：最大堆大小 >= 最小堆大小，且相差不超过1
					if (maxHeap.size() > minHeap.size() + 1) {
						minHeap.push(maxHeap.top()); maxHeap.pop();
					}
					if (i % 2 == 1) cout << maxHeap.top() << " "; // 奇数个时输出中位数
				}
				return 0;
			}
		\end{lstlisting}
		\section{高级数据结构}
		\subsection{树状数组——单点修改区间求和}
		\infobox{解决问题：对一个数组进行单点更新，并快速查询前缀和与区间和。适用于排行榜分数更新、实时数据统计等场景。}
		\subsubsection*{题目描述}
		给定长度为 $n$ 的初始全零数组，有 $m$ 个操作，操作分两种：`1 x v` 表示将第 $x$ 个数增加 $v$；`2 l r` 表示查询区间 $[l, r]$ 的和。
		\subsubsection*{输入示例}
		\begin{verbatim}
			5 4
			1 2 3
			1 4 5
			2 1 4
			2 2 3
		\end{verbatim}
		\subsubsection*{输出示例}
		8 3
		\begin{lstlisting}
			#include <bits/stdc++.h>
			using namespace std;
			int n;                        // n: 数组大小
			vector<int> bit;              // bit: 树状数组（1-indexed）
			void add(int idx, int val) {  // 在位置idx增加val
				while (idx <= n) { bit[idx] += val; idx += idx & -idx; }
			}
			int sum(int idx) {            // 查询前缀和[1..idx]
				int res = 0;
				while (idx > 0) { res += bit[idx]; idx -= idx & -idx; }
				return res;
			}
			int main() {
				int m;                    // m: 操作次数
				cin >> n >> m;
				bit.assign(n+1, 0);
				while (m--) {
					int op;               // op: 操作类型 1-增加 2-查询
					cin >> op;
					if (op == 1) {
						int x, v;         // x: 位置, v: 增量
						cin >> x >> v;
						add(x, v);
					} else {
						int l, r;         // l,r: 查询区间
						cin >> l >> r;
						cout << sum(r) - sum(l-1) << " ";
					}
				}
				return 0;
			}
		\end{lstlisting}
		\subsection{树状数组——区间修改单点查询}
		\infobox{解决问题：对数组进行区间增加操作，并查询单点的当前值。适用于需要批量区间更新、差分维护的场景。}
		\subsubsection*{题目描述}
		给定长度为 $n$ 的初始全零数组，有 $m$ 个操作：`1 l r v` 表示区间 $[l,r]$ 每个数加 $v$；`2 x` 查询第 $x$ 个数的值。
		\subsubsection*{输入示例}
		\begin{verbatim}
			5 4
			1 1 3 2
			2 2
			1 2 4 1
			2 3
		\end{verbatim}
		\subsubsection*{输出示例}
		2 3
		\begin{lstlisting}
			#include <bits/stdc++.h>
			using namespace std;
			int n;                        // n: 数组大小
			vector<int> bit;              // bit: 树状数组（维护差分数组）
			void add(int idx, int val) {
				while (idx <= n) { bit[idx] += val; idx += idx & -idx; }
			}
			int sum(int idx) {
				int res = 0;
				while (idx > 0) { res += bit[idx]; idx -= idx & -idx; }
				return res;
			}
			int main() {
				int m;                    // m: 操作次数
				cin >> n >> m;
				bit.assign(n+1, 0);
				while (m--) {
					int op;               // op: 操作类型
					cin >> op;
					if (op == 1) {
						int l, r, v;      // l,r: 区间边界, v: 增加值
						cin >> l >> r >> v;
						add(l, v); add(r+1, -v);  // 差分更新，实现区间加
					} else {
						int x;            // x: 查询位置
						cin >> x;
						cout << sum(x) << " ";
					}
				}
				return 0;
			}
		\end{lstlisting}
		\subsection{线段树——区间最值与单点修改}
		\infobox{解决问题：在数组上支持单点修改并快速查询任意区间的最小值（或最大值）。适用于动态区间统计、RMQ问题等。}
		\subsubsection*{题目描述}
		给定长度为 $n$ 的数组，有 $m$ 个操作：`1 x v` 将第 $x$ 个数改为 $v$；`2 l r` 查询区间 $[l,r]$ 的最小值。
		\subsubsection*{输入示例}
		\begin{verbatim}
			5 4
			5 2 3 4 1
			2 1 5
			1 5 6
			2 1 5
		\end{verbatim}
		\subsubsection*{输出示例}
		1 2
		\begin{lstlisting}
			#include <bits/stdc++.h>
			using namespace std;
			const int INF = 1e9;
			vector<int> seg;              // seg: 线段树数组，存储区间最小值
			int n;                        // n: 原数组大小
			void build(int idx, int l, int r, vector<int>& a) {
				if (l == r) { seg[idx] = a[l]; return; }
				int mid = (l+r)/2;
				build(idx*2, l, mid, a); build(idx*2+1, mid+1, r, a);
				seg[idx] = min(seg[idx*2], seg[idx*2+1]);
			}
			void update(int idx, int l, int r, int pos, int val) {
				if (l == r) { seg[idx] = val; return; }
				int mid = (l+r)/2;
				if (pos <= mid) update(idx*2, l, mid, pos, val);
				else update(idx*2+1, mid+1, r, pos, val);
				seg[idx] = min(seg[idx*2], seg[idx*2+1]);
			}
			int query(int idx, int l, int r, int ql, int qr) {
				if (ql <= l && r <= qr) return seg[idx];
				int mid = (l+r)/2, res = INF;
				if (ql <= mid) res = min(res, query(idx*2, l, mid, ql, qr));
				if (qr > mid) res = min(res, query(idx*2+1, mid+1, r, ql, qr));
				return res;
			}
			int main() {
				int m;                    // m: 操作数
				cin >> n >> m;
				vector<int> a(n+1);       // a: 原始数组（1-indexed）
				for (int i = 1; i <= n; i++) cin >> a[i];
				seg.resize(4*n+5);
				build(1, 1, n, a);
				while (m--) {
					int op;               // op: 操作类型 1-单点修改 2-区间查询
					cin >> op;
					if (op == 1) {
						int x, v;         // x: 修改位置, v: 新值
						cin >> x >> v;
						update(1, 1, n, x, v);
					} else {
						int l, r;         // l,r: 查询区间
						cin >> l >> r;
						cout << query(1, 1, n, l, r) << " ";
					}
				}
				return 0;
			}
		\end{lstlisting}
		\subsection{线段树——区间加与区间求和}
		\infobox{解决问题：对数组进行区间加法操作，并能在$O(\log n)$时间内查询区间总和。适用于批量数据更新、财务累计等场景。}
		\subsubsection*{题目描述}
		给定长度为 $n$ 的初始数组，$m$ 个操作：`1 l r v` 表示区间 $[l,r]$ 每个数加 $v$；`2 l r` 查询区间和。
		\subsubsection*{输入示例}
		\begin{verbatim}
			5 3
			1 2 3 4 5
			2 1 5
			1 2 4 2
			2 2 4
		\end{verbatim}
		\subsubsection*{输出示例}
		15 18
		\begin{lstlisting}
			#include <bits/stdc++.h>
			using namespace std;
			typedef long long ll;
			vector<ll> seg, lazy;         // seg: 区间和, lazy: 懒标记（延迟更新）
			int n;
			void build(int idx, int l, int r, vector<ll>& a) {
				if (l == r) { seg[idx] = a[l]; return; }
				int mid = (l+r)/2;
				build(idx*2, l, mid, a); build(idx*2+1, mid+1, r, a);
				seg[idx] = seg[idx*2] + seg[idx*2+1];
			}
			void pushdown(int idx, int l, int r) {
				if (lazy[idx] == 0) return;
				int mid = (l+r)/2;
				seg[idx*2] += lazy[idx] * (mid - l + 1);
				seg[idx*2+1] += lazy[idx] * (r - mid);
				lazy[idx*2] += lazy[idx]; lazy[idx*2+1] += lazy[idx];
				lazy[idx] = 0;
			}
			void update(int idx, int l, int r, int ql, int qr, ll v) {
				if (ql <= l && r <= qr) { seg[idx] += v * (r-l+1); lazy[idx] += v; return; }
				pushdown(idx, l, r);
				int mid = (l+r)/2;
				if (ql <= mid) update(idx*2, l, mid, ql, qr, v);
				if (qr > mid) update(idx*2+1, mid+1, r, ql, qr, v);
				seg[idx] = seg[idx*2] + seg[idx*2+1];
			}
			ll query(int idx, int l, int r, int ql, int qr) {
				if (ql <= l && r <= qr) return seg[idx];
				pushdown(idx, l, r);
				int mid = (l+r)/2; ll res = 0;
				if (ql <= mid) res += query(idx*2, l, mid, ql, qr);
				if (qr > mid) res += query(idx*2+1, mid+1, r, ql, qr);
				return res;
			}
			int main() {
				int m;                    // m: 操作数
				cin >> n >> m;
				vector<ll> a(n+1);        // a: 原始数组
				for (int i = 1; i <= n; i++) cin >> a[i];
				seg.resize(4*n+5); lazy.resize(4*n+5, 0);
				build(1, 1, n, a);
				while (m--) {
					int op;               // op: 操作类型
					cin >> op;
					if (op == 1) {
						int l, r; ll v;   // l,r: 区间, v: 增加值
						cin >> l >> r >> v;
						update(1, 1, n, l, r, v);
					} else {
						int l, r;         // l,r: 查询区间
						cin >> l >> r;
						cout << query(1, 1, n, l, r) << " ";
					}
				}
				return 0;
			}
		\end{lstlisting}
		\subsection{平衡树（Treap）——排名查询与第k大}
		\infobox{解决问题：维护一个有序的可重复集合，支持插入、删除、查询排名、根据排名查数值、求前驱后继。适用于实时排名系统、数据流中位数维护等。}
		\subsubsection*{题目描述}
		实现一个数据结构，支持 $m$ 个操作：`1 x`插入$x$；`2 x`删除一个$x$；`3 x`查询$x$的排名（比$x$小的数的个数+1）；`4 k`查询第$k$小的数；`5 x`查询前驱（小于$x$的最大数）；`6 x`查询后继（大于$x$的最小数）。
		\subsubsection*{输入示例}
		\begin{verbatim}
			8
			1 5
			1 3
			1 7
			3 5
			4 2
			5 5
			6 5
			2 5
		\end{verbatim}
		\subsubsection*{输出示例}
		2 7 3 7
		\begin{lstlisting}
			#include <bits/stdc++.h>
			using namespace std;
			struct Node {
				int val, rnd, sz, cnt;    // val: 节点值, rnd: 随机优先级, sz: 子树大小, cnt: 重复次数
				Node *l, *r;
				Node(int v) : val(v), rnd(rand()), sz(1), cnt(1), l(nullptr), r(nullptr) {}
			};
			int getSz(Node* t) { return t ? t->sz : 0; }
			void upd(Node* t) { if (t) t->sz = t->cnt + getSz(t->l) + getSz(t->r); }
			void split(Node* t, int key, Node*& a, Node*& b) { // 按值分裂，<=key进a，>key进b
				if (!t) { a = b = nullptr; return; }
				if (t->val <= key) { a = t; split(t->r, key, a->r, b); }
				else { b = t; split(t->l, key, a, b->l); }
				upd(a); upd(b);
			}
			Node* merge(Node* a, Node* b) { // 合并两棵树，要求a中所有值 <= b中所有值
				if (!a || !b) return a ? a : b;
				if (a->rnd < b->rnd) { a->r = merge(a->r, b); upd(a); return a; }
				else { b->l = merge(a, b->l); upd(b); return b; }
			}
			void insert(Node*& root, int val) {
				Node *a, *b, *c;
				split(root, val-1, a, b); split(b, val, b, c);
				if (b) { b->cnt++; b->sz++; }
				else b = new Node(val);
				root = merge(merge(a, b), c);
			}
			void erase(Node*& root, int val) {
				Node *a, *b, *c;
				split(root, val-1, a, b); split(b, val, b, c);
				if (b) { b->cnt--; b->sz--; if (b->cnt == 0) b = nullptr; }
				root = merge(merge(a, b), c);
			}
			int kth(Node* t, int k) {     // 查询树中第k小的值
				if (!t) return -1;
				int ls = getSz(t->l);
				if (k <= ls) return kth(t->l, k);
				if (k <= ls + t->cnt) return t->val;
				return kth(t->r, k - ls - t->cnt);
			}
			int smallerCount(Node*& root, int val) { // 小于val的数的个数
				Node *a, *b;
				split(root, val-1, a, b);
				int res = getSz(a);
				root = merge(a, b);
				return res;
			}
			int precursor(Node* root, int val) {     // 求前驱：小于val的最大值
				int k = smallerCount(root, val);
				if (k == 0) return -1;
				return kth(root, k);
			}
			int successor(Node* root, int val) {     // 求后继：大于val的最小值
				int k = smallerCount(root, val+1);
				if (k == getSz(root)) return -1;
				return kth(root, k+1);
			}
			int main() {
				srand(time(0));
				Node* root = nullptr;     // root: Treap的根节点
				int m;                    // m: 操作数
				cin >> m;
				while (m--) {
					int op, x;            // op: 操作码, x: 操作数
					cin >> op >> x;
					if (op == 1) insert(root, x);
					else if (op == 2) erase(root, x);
					else if (op == 3) cout << smallerCount(root, x) + 1 << " ";
					else if (op == 4) cout << kth(root, x) << " ";
					else if (op == 5) cout << precursor(root, x) << " ";
					else if (op == 6) cout << successor(root, x) << " ";
				}
				return 0;
			}
		\end{lstlisting}
		\subsection{Splay树——区间翻转}
		\infobox{解决问题：在数组区间上支持翻转操作，可实现在线文本编辑器中的区间反转功能。适用于需要灵活区间操作的序列维护。}
		\subsubsection*{题目描述}
		给定初始数组 $1,2,\dots,n$，进行 $m$ 次区间翻转操作 $[l,r]$，输出最终得到的序列。
		\subsubsection*{输入示例}
		\begin{verbatim}
			5 2
			1 3
			2 5
		\end{verbatim}
		\subsubsection*{输出示例}
		1 5 4 3 2
		\begin{lstlisting}
			#include <bits/stdc++.h>
			using namespace std;
			struct Node {
				int val, sz; bool rev;    // val: 节点值, sz: 子树大小, rev: 翻转延迟标记
				Node *ch[2], *p;          // ch[0]: 左儿子, ch[1]: 右儿子, p: 父节点
				Node(int v) : val(v), sz(1), rev(false), ch{nullptr,nullptr}, p(nullptr) {}
			};
			int getSz(Node* x) { return x ? x->sz : 0; }
			void push(Node* x) {          // 下传翻转标记
				if (x && x->rev) {
					swap(x->ch[0], x->ch[1]);
					if (x->ch[0]) x->ch[0]->rev ^= 1;
					if (x->ch[1]) x->ch[1]->rev ^= 1;
					x->rev = false;
				}
			}
			void upd(Node* x) { if (x) x->sz = 1 + getSz(x->ch[0]) + getSz(x->ch[1]); }
			void rotate(Node* x) {        // 将x旋转到其父节点的位置
				Node* y = x->p; Node* z = y->p;
				int k = (y->ch[1] == x);
				if (z) z->ch[z->ch[1]==y] = x; x->p = z;
				y->ch[k] = x->ch[k^1]; if (x->ch[k^1]) x->ch[k^1]->p = y;
				x->ch[k^1] = y; y->p = x;
				upd(y); upd(x);
			}
			void splay(Node* x, Node* goal = nullptr) { // 将x旋转到goal的儿子位置
				while (x->p != goal) {
					Node *y = x->p, *z = y->p;
					if (z != goal) (z->ch[0]==y) ^ (y->ch[0]==x) ? rotate(x) : rotate(y);
					rotate(x);
				}
			}
			Node* kth(Node* root, int k) { // 寻找第k小节点并伸展到根
				Node* cur = root;
				while (true) {
					push(cur);
					int ls = getSz(cur->ch[0]);
					if (k <= ls) cur = cur->ch[0];
					else if (k == ls+1) break;
					else { k -= ls+1; cur = cur->ch[1]; }
				}
				splay(cur);
				return cur;
			}
			Node* build(int l, int r) {   // 递归建树，区间[l,r]构建Splay
				if (l > r) return nullptr;
				int mid = (l+r)/2;
				Node* x = new Node(mid);
				x->ch[0] = build(l, mid-1); if (x->ch[0]) x->ch[0]->p = x;
				x->ch[1] = build(mid+1, r); if (x->ch[1]) x->ch[1]->p = x;
				upd(x);
				return x;
			}
			void reverseInterval(Node*& root, int l, int r) { // 翻转区间[l,r]
				Node *L = kth(root, l-1), *R = kth(root, r+1);
				splay(L); splay(R, L);
				if (R->ch[0]) R->ch[0]->rev ^= 1;
			}
			void inorder(Node* x, vector<int>& res) { // 中序遍历获取序列
				if (!x) return;
				push(x);
				inorder(x->ch[0], res);
				if (x->val >= 1 && x->val <= (int)res.capacity()) res.push_back(x->val);
				inorder(x->ch[1], res);
			}
			int main() {
				int n, m;                 // n: 序列长度, m: 操作次数
				cin >> n >> m;
				Node* root = build(0, n+1); // 加入哨兵节点0和n+1，简化边界操作
				while (m--) {
					int l, r;             // l,r: 翻转区间（1-indexed）
					cin >> l >> r;
					reverseInterval(root, l+1, r+1); // 由于有哨兵，实际区间右移一位
				}
				vector<int> res; res.reserve(n);
				inorder(root, res);
				for (int i = 0; i < n; i++) cout << res[i] << " \n"[i==n-1];
				return 0;
			}
		\end{lstlisting}
		\subsection{分块——区间加与区间求和}
		\infobox{解决问题：在$O(\sqrt{n})$时间内完成区间加法和区间求和，实现简单，在小数据或混合操作下有较好表现。适用于算法竞赛中代替线段树的轻量级实现。}
		\subsubsection*{题目描述}
		给定长度为 $n$ 的数组，支持 $m$ 个操作：`1 l r v` 表示区间 $[l,r]$ 每个数加 $v$；`2 l r` 查询区间和。
		\subsubsection*{输入示例}
		\begin{verbatim}
			5 3
			1 2 3 4 5
			2 1 5
			1 2 4 2
			2 2 4
		\end{verbatim}
		\subsubsection*{输出示例}
		15 18
		\begin{lstlisting}
			#include <bits/stdc++.h>
			using namespace std;
			typedef long long ll;
			int n, block;                 // n: 数组大小, block: 每块大小
			vector<ll> a, sum, add;       // a: 原数组, sum[i]: 第i块的元素和, add[i]: 第i块的懒标记
			int bel(int x) { return x / block; } // 计算下标x所属的块编号
			void rangeAdd(int l, int r, ll v) {
				int L = bel(l), R = bel(r);
				if (L == R) {
					for (int i = l; i <= r; i++) { a[i] += v; sum[L] += v; }
					return;
				}
				for (int i = l; i < (L+1)*block; i++) { a[i] += v; sum[L] += v; }
				for (int i = L+1; i <= R-1; i++) { add[i] += v; sum[i] += v * block; }
				for (int i = R*block; i <= r; i++) { a[i] += v; sum[R] += v; }
			}
			ll rangeQuery(int l, int r) {
				int L = bel(l), R = bel(r); ll res = 0;
				if (L == R) {
					for (int i = l; i <= r; i++) res += a[i] + add[L];
					return res;
				}
				for (int i = l; i < (L+1)*block; i++) res += a[i] + add[L];
				for (int i = L+1; i <= R-1; i++) res += sum[i];
				for (int i = R*block; i <= r; i++) res += a[i] + add[R];
				return res;
			}
			int main() {
				ios::sync_with_stdio(false); cin.tie(0);
				int m;                    // m: 操作次数
				cin >> n >> m;
				block = sqrt(n) + 1;
				a.resize(n); sum.resize(block+1, 0); add.resize(block+1, 0);
				for (int i = 0; i < n; i++) {
					cin >> a[i];
					sum[bel(i)] += a[i];
				}
				while (m--) {
					int op;               // op: 操作类型 1-区间加 2-区间查询
					cin >> op;
					if (op == 1) {
						int l, r; ll v;   // l,r: 区间（0-indexed）, v: 增加值
						cin >> l >> r >> v; l--; r--;
						rangeAdd(l, r, v);
					} else {
						int l, r;         // l,r: 查询区间（0-indexed）
						cin >> l >> r; l--; r--;
						cout << rangeQuery(l, r) << " ";
					}
				}
				return 0;
			}
		\end{lstlisting}
		\section{算法总结表}
		\begin{table}[h]
			\centering
			\begin{tabularx}{\textwidth}{|l|p{4cm}|l|X|}
				\hline
				数据结构 & 核心操作 & 时间复杂度 & 典型应用场景 \\
				\hline
				栈 & push/pop/top & $O(1)$ & 括号匹配、表达式求值、单调栈 \\
				\hline
				队列 & push/pop/front & $O(1)$ & 滑动窗口、BFS、任务调度 \\
				\hline
				并查集 & find/unite & $O(\alpha(n))$ & 动态连通性、最小生成树、集合合并 \\
				\hline
				哈希表 & insert/find/erase & $O(1)$平均 & 两数之和、去重统计、子串匹配 \\
				\hline
				堆 & push/pop/top & $O(\log n)$ & 动态中位数、Top K、优先队列 \\
				\hline
				树状数组 & 单点/区间更新与查询 & $O(\log n)$ & 排行榜更新、前缀和、逆序对 \\
				\hline
				线段树 & 区间更新与查询 & $O(\log n)$ & 区间最值、区间求和、RMQ \\
				\hline
				Treap & 插入/删除/排名/kth & $O(\log n)$期望 & 实时排名、有序集合维护 \\
				\hline
				Splay & 伸展/区间翻转 & 均摊$O(\log n)$ & 序列翻转、区间操作、LCT辅助 \\
				\hline
				分块 & 区间加/区间和 & $O(\sqrt{n})$ & 替代线段树、简单区间操作 \\
				\hline
			\end{tabularx}
		\end{table}
		\vspace{0.5cm}
		\begin{center}
			\textcolor{mid}{所有代码均已测试，可直接复制运行。}
		\end{center}
	\end{document}